% Generated from verified GitHub CSV at 7e136fc. See evidence/SOURCE_MANIFEST.json.
\begin{table*}[t]
\centering\footnotesize
\caption{기존 직사각형의 재질별 보정 효과. 각 재질 안에서 같은 영상·참조의 전후를 비교한다.}
\label{sup:type_results}
\setlength{\tabcolsep}{3pt}
\begin{tabularx}{\textwidth}{Ylrrrrrrr}
\toprule
집단 & 방법 & 영상 & \shortstack{코너\\ 참조/유효 예측} & 중앙값(px) & P90(px) & \pckten(\%) & 이동(cm) & 회전(도) \\
\midrule
직사각형 / 플라스틱 & Base & 194 & 1513/1459 & 7.509 & 43.637 & 60.079 & 10.472 & 2.450 \\
 & P & 194 & 1513/1459 & 6.397 & 42.254 & 64.596 & 9.703 & 2.099 \\
 & N2 & 194 & 1513/1459 & 6.216 & 41.995 & 65.719 & 9.259 & 2.045 \\
 & N3 & 194 & 1513/1459 & 6.203 & 41.703 & 65.675 & 9.363 & 2.017 \\
\midrule
직사각형 / 목재 & Base & 125 & 986/986 & 6.124 & 45.540 & 68.560 & 4.204 & 2.650 \\
 & P & 125 & 986/986 & 5.371 & 45.850 & 71.941 & 3.850 & 2.331 \\
 & N2 & 125 & 986/986 & 5.243 & 44.473 & 72.989 & 3.805 & 2.194 \\
 & N3 & 125 & 986/986 & 5.274 & 44.865 & 73.056 & 3.739 & 2.202 \\
\bottomrule
\end{tabularx}
\tabnote{플라스틱 194장과 목재 125장의 합은 319장이다. 중앙 코너 오차·P90은 매칭된 유효 코너, \pckten은 전체 참조 코너 분모를 사용한다. 이동·회전은 중앙값이다. 서로 다른 재질의 난도를 동일하다고 가정하지 않는다.}
\end{table*}
